\noindent\textbf{Resultado.} $y(x)=\tfrac13e^{-x}+e^{-4x}$, $x\in\mathbb R$.

\noindent\textit{Desarrollo.} El factor integrante es $e^{4x}$, de modo que $(e^{4x}y)'=e^{3x}$. Integrar da $y=\frac13e^{-x}+Ce^{-4x}$. El dato inicial determina $C=1$. Al derivar y sumar $4y$, el término $e^{-4x}$ se cancela y queda $e^{-x}$; en cero la suma es $4/3$. Los coeficientes continuos garantizan unicidad en toda la recta.
